% NOTE: template from https://www.maltegerken.de/blog/2021/07/creating-pdf-forms-with-latex/
\documentclass[parskip=full]{scrartcl}

\usepackage{tikz}
\usepackage{hyperref}
\usepackage{setspace}
\usepackage{pbox}

\renewcommand*{\DefaultOptionsofRadio}{print,radio, radiosymbol=6, width=\baselineskip, bordercolor={black}, borderwidth=0pt}
\renewcommand*{\DefaultOptionsofText}{print,bordercolor={black}, backgroundcolor=white, borderwidth=0pt}

\renewcommand{\LayoutTextField}[2]{% label, field
\setbox0=\hbox{#1\unskip}\ifdim\wd0=0pt
\setbox1=\hbox{#2\unskip}\ifdim\ht1>3ex
% Multiline
\begin{tikzpicture}[every node/.style={inner sep=0,outer sep=0}]
  \node[anchor=west] (TextFieldNode) at (0cm,0cm) {#2};
  \draw [thick] (current bounding box.south west) rectangle (current bounding box.north east);
\end{tikzpicture}%
\else
% Inline field, lowered a little bit to be better integrated into the text
\raisebox{-3.2pt}{\begin{tikzpicture}[every node/.style={inner sep=0,outer sep=0}]
  \node[anchor=west] (TextFieldNode) at (0cm,0cm) {#2};
  \draw[thick] ([yshift=-0.3ex]TextFieldNode.south west) -- ([yshift=-0.3ex]TextFieldNode.south east);
\end{tikzpicture}}%
\fi
\else
% Field with label below it
\begin{tikzpicture}[every node/.style={inner sep=0,outer sep=0}]
  \node[anchor=west] (TextFieldNode) at (0cm,2ex) {#2};
  \draw[thick] ([yshift=-0.3ex]TextFieldNode.south west) -- ([yshift=-0.3ex]TextFieldNode.south east);
  \node[anchor=west,font=\footnotesize] at (0cm,-0.9ex) {#1};
\end{tikzpicture}%
\fi
}
\newcommand{\field}[2]{\TextField[width=#2]{#1}}
\newcommand{\fieldinline}[2]{\TextField[width=#2,name=#1]{}}

\newcommand{\radiosize}{0.33cm}
\newcommand{\yesnoticks}[1]{%
\raisebox{-3.2pt}{\begin{tikzpicture}[every node/.style={inner sep=0,outer sep=0}]
  \node[anchor=west,style={inner sep=2px}] (FieldYes) at (0cm,0cm) {\ChoiceMenu[radio=true,name=#1,width=\radiosize,height=\radiosize]{}{=Yes}};
  \node[anchor=west] (LabelYes) at ([xshift=0.7ex]FieldYes.east) {Yes};
  \node[anchor=west,style={inner sep=2px}] (FieldNo) at ([xshift=1ex]LabelYes.east) {\ChoiceMenu[radio=true,name=#1,width=\radiosize,height=\radiosize]{}{=No}};
  \node[anchor=west] (LabelNo) at ([xshift=0.7ex]FieldNo.east) {No};
  \draw [thick] ([xshift=-(\radiosize+0.15cm),yshift=-(\radiosize+0.15cm)]FieldYes.north east) rectangle (FieldYes.north east);
  \draw [thick] ([xshift=-(\radiosize+0.15cm),yshift=-(\radiosize+0.15cm)]FieldNo.north east) rectangle (FieldNo.north east);
\end{tikzpicture}}%
}
\newcommand{\yesno}[2]{\pbox{0.7\textwidth}{\setstretch{1}#1}\hfill\yesnoticks{#2}}

\begin{document}
  \section*{Glasgow Hackerspace CIC membership form}
  \begin{Form}
    \field{email}{1.0\linewidth}%

    \field{discord username (if you are on the discord)}{1.0\linewidth}%

    \field{pseudo or name you want to be known by}{1.0\linewidth}%

    \field{Membership fee per month (see below for more information)}{1.0\linewidth}%

    \yesno{Are you over 18 years old?}{adult}

    \subsection*{Postal contact informations}

    \field{first name}{0.45\linewidth}\hfill\field{last name}{0.5\linewidth}%

    \field{address}{1.0\linewidth}%

    \field{city}{0.3\linewidth}\hfill\field{postcode}{0.3\linewidth}\hfill\field{country}{0.3\linewidth}%
  \end{Form}

  We will contact you by email as soon as possible to confirm your membership.

  Suggested membership is £25/month. Membership fees are currently optional, until we find a permanent space, but any membership fees will go towards securing a space, building up a buffer for rent, and possibly purchasing equipment for the space.

  Our Articles of Association are available at: https://github.com/glasgowhackerspace/company

  \textbf{Disclaimer:} Your informations will be stored for 10 years after the end of your membership, in line with UK company law. (section 121 of Companies Act 2006)
  They will be stored in Proton Drive and accessible only to the company directors; or per reasonable request to inspect. (section 116 of Companies Act 2006)
\end{document}
